% !PRINTED 519
\weberpage{519}
%<<CONT
then each $\varphi(\alpha_i, y_\alpha)$ must be equal to zero, that is, $\varphi(x, y_\alpha)$ must
be divisible by $f(x, y_\alpha)$. But it does not follow from this that
$f(x, y_\alpha) = f_\alpha(x, \psi)$ is also irreducible in the field $\Omega(\psi)$.

These are the same theorems as in \S.\,151. It is further shown
there that for $y_\alpha$ one can take a function of $\alpha$ alone,
which we shall not go into again.

It is of interest, however, to follow the connection between the \gesperrt{Galois}
groups of the individual factors $f(x, y_\alpha), f(x, y_\beta), \ldots f(x, y_\sigma)$.
We have already seen above that one obtains the
group of the factor $f(x, y_\alpha)$ in the field $\Omega(\psi)$ by seeking out
all the permutations of the $\alpha$ that are produced by $Q$.
We shall denote this group by $P_\alpha$, and likewise
use the signs $P_\beta, \ldots P_\sigma$ for the other factors.

We select from $P$ any $s - 1$ permutations $\pi_\beta \ldots \pi_\sigma$
of such a kind that by $\pi_\beta$ the row $A$ passes into $B$ etc.,
by $\pi_\sigma$ the row $A$ into $S$; and for the sake of
simplicity we shall now — which is evidently still entirely open to us — assume the
designation of the $\beta, \ldots \sigma$ so chosen that $\pi_\beta, \ldots \pi_\sigma$
produce the following permutations:
\[
\begin{pmatrix} \alpha_0, & \alpha_1 & \ldots & \alpha_{r-1} \\ \beta_0, & \beta_1 & \ldots & \beta_{r-1} \end{pmatrix}, \ldots
\begin{pmatrix} \alpha_0, & \alpha_1 & \ldots & \alpha_{r-1} \\ \sigma_0, & \sigma_1 & \ldots & \sigma_{r-1} \end{pmatrix}.
\]
Now let $\varkappa$ be any element of $Q$ which among the $\alpha$ effects the
permutation
\[
(\alpha) = \begin{pmatrix} \alpha_0, & \alpha_1 & \ldots & \alpha_{r-1} \\ \alpha_{a_0}, & \alpha_{a_1} & \ldots & \alpha_{a_{r-1}} \end{pmatrix}
\]
then
\[
\varkappa_\beta = \pi_\beta^{-1} \varkappa \pi_\beta
\]
is likewise contained in $Q$, and by $\varkappa_\beta$ there is produced among the $\beta$ the
permutation
\[
(\beta) = \begin{pmatrix} \beta_0, & \beta_1 & \ldots & \beta_{r-1} \\ \beta_{a_0}, & \beta_{a_1} & \ldots & \beta_{a_{r-1}} \end{pmatrix},
\]
and hence among the indices of the $\beta$ exactly the same permutation
as by $\varkappa$ among the indices of $\alpha$. If we take
$(\beta)$ and a corresponding $\varkappa_\beta$ as given, then $\pi_\beta \varkappa_\beta \pi_\beta^{-1}$
yields $\varkappa$ again, so that the permutations of the $\beta$
produced by $Q$ agree exactly with those which are produced among the
$\alpha$; and the same holds for the remaining
rows. But with this the theorem is proved.
